\title{%
  Models in security
}
\author{Daniel Bosk\thanks{%
    This material is available under the Creative Commons 
    Attribution-NonCommercial-ShareAlike (CC-BY-NC-SA) 4.0 international 
    license.
    The material was written with some aid from GitHub Copilot.
}}
\institute{%
  KTH EECS
}

\begin{frame}
  \maketitle
\end{frame}

\mode*

\begin{abstract}
  \input{abstract.tex}
\end{abstract}

\clearpage

\section{Part I}

\begin{frame}
  \begin{block}{From your Part I reflections: recurring themes}
    \mode<presentation>{%
      \begin{itemize}
        \item Proofs hold in idealised models; practice breaks assumptions 
          (side channels, timing, EM)
        \item Models are still essential: give baselines, definitions, and 
          comparability
        \item Security is layered: crypto, implementation, hardware, human 
          factors
        \item TCB perspective: failures outside the model can break the system
        \item Assumptions matter: clarify adversary capabilities and scope
      \end{itemize}
    }
    \mode<article>{%
      Many of you stressed that cryptographic proofs establish guarantees only 
      within explicit, idealised models, while practical systems violate 
      assumptions via side channels (EM/timing/power), protocol metadata, and 
      human factors. Yet, models remain indispensable because they define 
      security precisely, provide a baseline for comparison, and help bound 
      claims. Several reflections framed this in TCB terms: if any component 
      outside the model fails, the whole may fail. Across responses, the call 
      was to make assumptions and attacker capabilities explicit and to combine 
      formal models with empirical testing.
    }
  \end{block}
\end{frame}

\begin{frame}
  \begin{block}{Part I — What you said about provable security vs reality}
  \mode<presentation>{
    \begin{description}
      \item[Consensus] Proofs give baseline guarantees under explicit 
        assumptions; implementations and environments add side channels and 
        leakage
      \item[Tensions] What belongs inside the model vs out-of-scope? 
        Definitions of secure differ (cryptographic vs empirical)
      \item[Value] Models enable comparison, clarity, and rejection of designs 
        that fail even under idealized conditions
    \end{description}
  }
  \mode<article>{
    Many students stressed that proofs certify properties under a precise 
    threat model and assumptions, while practice adds leakage (timing, EM, 
    caching, human factors). Models remain useful as baselines for design and 
    comparison, but they do not close the gap to reality; testing and 
    engineering must complement them. Several noted the Trusted Computing Base 
    perspective: a flaw outside the model can break the whole system.
  }
  \end{block}
\end{frame}

\begin{frame}
  \begin{question}[Discussion (Part I)]
    \begin{enumerate}
      \item Where should we draw the threat boundary in a proof? Should side 
        channels be inside or outside the cryptographic model?
      \item When is a proof \emph{misleading} rather than merely 
        \emph{limited}? What signals that a model is too narrow?
      \item How do we use models to prioritise mitigations: algorithm redesign 
        vs. implementation hardening vs. operational controls?
      \item If a scheme fails in practice but not in theory, what should change 
        first: the model, the implementation, or the deployment context?
    \end{enumerate}
  \end{question}
\end{frame}

\begin{frame}
  \begin{activity}[Discuss — Where should the boundary of the model be?]
  \begin{itemize}
    \item For each leakage below, decide how to deal with it scientifically:
      \begin{itemize}
        \item Packet length and inter-keystroke timing in SSH
        \item EM emanations from windowed exponentiation
      \end{itemize}
    \item If you move one item into the model, what new assumption or cost 
      appears?
  \end{itemize}
  \mode<article>{
    Prompt students to justify modeling choices explicitly (e.g., include 
    timing in the protocol model implies constant-time framing or padding, with 
    performance cost).
  }
  \end{activity}
\end{frame}

% XXX revise below

\begin{frame}
  \begin{activity}[Activity — Diagnosing a proof-to-practice gap]
  \begin{itemize}
    \item Pick a scheme with a reductionist proof (e.g., ElGamal under DDH)
    \item Identify:
      \begin{description}
        \item[Claim] What the proof guarantees
        \item[Assumptions] Hardness and adversary capabilities
        \item[Caveat] One realistic behavior the proof omits
        \item[Mitigation] A protocol or implementation defense to restore the guarantee
      \end{description}
  \end{itemize}
  \mode<article>{
    This ties student reflections (baseline guarantees, omitted channels) to 
    actionable mitigations (padding, blinding, constant-time, shielding).
  }
  \end{activity}
\end{frame}


\section{Part II}

\begin{frame}
  \begin{block}{From your Part II reflections: what formal methods add}
    \mode<presentation>{%
      \begin{itemize}
        \item Rigor and reproducibility: machine-checked arguments, counterexamples
        \item Complementarity: theorem proving for specs; model checking for behaviours
        \item Limits: state explosion, finite-state modelling, cost and expertise
        \item Bridge to crypto: reductions define goals; formal methods check protocols/implementations
      \end{itemize}
    }
    \mode<article>{%
      You highlighted that formal verification and model checking bring deductive rigor and repeatability, offering machine-checked proofs and counterexamples that sharpen claims. Model checking excels for finite-state behaviours; theorem proving for general logical properties. Cost, expertise, and state explosion remain practical limits. Many linked this to cryptography: reductions justify a design's goal (\enquote{if DDH is hard, ElGamal is IND-CPA}), while formal methods verify protocols and implementations respect those goals (e.g., no secret-dependent control flow, correct key handling).
    }
  \end{block}
\end{frame}

\begin{frame}
  \begin{block}{Part II — What you said about formal methods}
  \mode<presentation>{
    \begin{description}
      \item[Formal verification] Deductive, math-based proofs for code/specs; high assurance, but effortful and assumption-bound
      \item[Model checking] Exhaustive exploration of finite models; counterexamples as scientific evidence; faces state explosion
      \item[Relation to crypto] Proofs link security to assumptions; formal methods check protocols/implementations meet those specs
    \end{description}
  }
  \mode<article>{
    Students highlighted rigor, reproducibility, and clear assumptions. Many positioned crypto proofs for primitives, model checking for protocols, and formal verification for implementations, noting remaining limits (side channels, model scope).
  }
  \end{block}
\end{frame}

\begin{frame}
  \begin{question}[Discussion (Part II)]
    \begin{enumerate}
      \item What should we formally verify for a cryptosystem stack: primitives, protocols, code, or all of the above?
      \item Where is model checking appropriate in security (finite-state protocols, access control, APIs)? Where is it not?
      \item How far can formal methods go on side channels (e.g., constant-time proofs, leakage models) before costs outweigh benefits?
      %\item What guarantees do decision-makers actually need: safety 
        %properties, non-interference, or quantified risk bounds?
    \end{enumerate}
  \end{question}
\end{frame}

%\begin{frame}
%  \begin{activity}[Design exercise — What to prove, what to check]
%  \begin{itemize}
%    \item You are shipping a new crypto library:
%      \begin{description}
%        \item[Prove] One property you will verify with a theorem prover (e.g., 
%          IND-CPA of a mode under stated assumptions)
%        \item[Check] One protocol property you will model-check (e.g., no replay or downgrade)
%        \item[Enforce] One implementation property you will enforce (e.g., constant-time, nonce uniqueness)
%      \end{description}
%    \item Be explicit about the assumptions and the TCB at each layer.
%  \end{itemize}
%  \mode<article>{
%    Expect answers mapping layers to methods and surfacing assumptions (PRF security, reliable RNG, constant-time discipline).
%  }
%  \end{activity}
%\end{frame}

\begin{frame}
  \begin{activity}[Discuss — Limits of formal methods]
  \begin{itemize}
    \item When does adding a side-channel model to verification help, and when does it harm tractability?
    \item How would you validate that verified code remains secure on real hardware?
    \item What evidence will you ship with a system: proof artifacts, model-checker traces, tests?
  \end{itemize}
  \end{activity}
\end{frame}



\section{Part III}

%\begin{frame}
%  \begin{block}{Part III — What you said about Tor papers}
%  \mode<presentation>{
%    \begin{description}
%      \item[TorPS/UGR] Realistic routing + user models; good coverage, upper-bound risks; sensitivity to routing inference and behavior models
%      \item[Shadow] Real Tor code in simulated network; repeatable, comparable; cannot capture full Internet variability
%      \item[WF+WO] Stronger adversary via auxiliary info; clarifies upper bounds; oracle realism is key
%      \item[Online WF] Genuine Tor traffic; high external validity; ethics and reproducibility constraints; accuracy drops for larger monitored sets
%    \end{description}
%  }
%  \mode<article>{
%    Students contrasted isolating vs mirror models and stressed external validity, ethics, and the base-rate problem in WF results.
%  }
%  \end{block}
%\end{frame}

\begin{frame}
  \begin{block}{How the papers model Tor (from your Part III reflections)}
    \begin{description}
      \item[Users Get Routed] TorPS path simulation with AS/IXP adversaries and user behaviour models to estimate correlation risk over time
      \item[Shadow] Discrete-event simulation that runs unmodified Tor code; controlled, reproducible experiments at scale
      \item[WF + Website Oracles] Oracle-augmented adversary model combining WF with side information to reduce false positives
      \item[Online WF (real Tor)] Models adversary at exit (and opt-in entry) with genuine Tor traffic and online learning in open-world settings
    \end{description}
  \end{block}
\end{frame}

\begin{frame}
  \begin{block}{Methodological strengths vs weaknesses (as you noted)}
    \mode<presentation>{%
      \begin{description}
        \item[Strengths]
          \begin{itemize}
            \item Realism via live data (Online WF); repeatability and control (Shadow)
            \item Broad what-if exploration (TorPS); clarifies attacker knowledge (WO)
            \item Counterexamples and stress tests without harming users
          \end{itemize}
        \item[Weaknesses]
          \begin{itemize}
            \item Model error and idealisation (routing inference, perfect oracles)
            \item External validity gaps (synthetic crawls vs live heterogeneity)
            \item Ethical limits and reproducibility constraints with real users
          \end{itemize}
      \end{description}
    }
    \mode<article>{%
      Strengths included realism from live exit-relay data (Online WF), repeatability and control by running real Tor in simulation (Shadow), and the ability to probe broad adversary conditions (TorPS) and attacker side information (WO). Weaknesses centered on idealisations (e.g., routing inference errors, perfect oracles), synthetic-data optimism, and ethical constraints that both increased realism and reduced reproducibility for live evaluations.
    }
  \end{block}
\end{frame}

%\begin{frame}
%  \begin{block}{Consequences for results (synthesis of your points)}
%    \mode<presentation>{%
%      \begin{itemize}
%        \item Lab/synthetic WF often overestimates attack power; live Tor shows strong drop-offs beyond small monitored sets
%        \item WO-style side information can overturn \enquote{defences} by collapsing false positives
%        \item Simulators are excellent for design comparison and safety, but not substitutes for measurement
%        \item Attacker models drive conclusions: be explicit about capabilities, scope, and data sources
%      \end{itemize}
%    }
%    \mode<article>{%
%      A consistent pattern emerged: closed-world, synthetic WF yields optimistic accuracy that erodes under live, noisy, evolving Tor traffic, especially beyond small monitored site sets. Oracle-style side information can substantially reduce false positives, challenging prior defences. Simulators (Shadow) are invaluable for safe, comparative experiments but must be paired with measurement to gauge external validity. Overall, attacker capability assumptions dominate conclusions and must be made explicit.
%    }
%  \end{block}
%\end{frame}

\begin{frame}
  \begin{activity}[Small-group critique — Model virtues and validity]
  \begin{itemize}
    \item Papers: TorPS, Shadow, WF+WO, Online WF
    \item Fill in for each paper:
      \begin{description}
        \item[Model type] Isolating vs mirror; what is abstracted
        \item[Virtues] Similarity, robustness, precision, transparency, tractability
        \item[External validity] What claims generalize; what do not
        \item[Improvement] One concrete way to strengthen the methodology ethically
      \end{description}
  \end{itemize}
  \end{activity}
\end{frame}

%\begin{frame}
%  \begin{activity}[Quantitative reasoning — The base-rate trap in WF]
%  \begin{itemize}
%    \item Suppose an attacker monitors 5 vs 1000 sites; discuss how precision and recall change
%    \item With a website oracle, which metric becomes critical for users’ risk?
%    \item Propose an evaluation report that avoids misleading accuracy and addresses base rates
%  \end{itemize}
%  \mode<article>{
%    Drive toward precision at low prevalence, calibrated thresholds, and reporting PPV/NPV alongside recall.
%  }
%  \end{activity}
%\end{frame}

%\begin{frame}
%  \begin{activity}[Ethics and reproducibility]
%  \begin{itemize}
%    \item For studies on live Tor, list practices that balance validity and ethics:
%      \begin{itemize}
%        \item Opt-in circuits, anonymized labels, ephemeral keys, aggregate statistics
%        \item Release code, configs, and synthetic replayers seeded by public aggregates
%      \end{itemize}
%    \item What could each paper share to maximize reproducibility without risking users?
%  \end{itemize}
%  \end{activity}
%\end{frame}

%\begin{frame}
%  \begin{question}[Discussion (Part III)]
%    \begin{enumerate}
%      \item How realistic are proposed oracles in your threat models (DNS, CDN logs, RTB, resolvers)? What access do real adversaries have?
%      \item When is Shadow \emph{enough} to evaluate a change? What requires live measurement despite ethical/repro concerns?
%      \item Which metrics best capture external validity for Tor studies (concept drift, open-world, base rates, path churn)?
%      \item What defences target the \enquote{right} layer given these results: path selection, padding, traffic shaping, resolver policy, user guidance?
%    \end{enumerate}
%  \end{question}
%\end{frame}

% ---------------- Wrap-up ----------------

\section{Wrap-up}

\begin{frame}
  \begin{block}{Takeaways across Parts I–III}
    \begin{itemize}
      \item Be explicit: assumptions, attacker capabilities, and the model’s scope
      \item Combine methods: reductions + formal verification + simulation + measurement
      \item Judge validity: internal (control), construct (right properties), external (realism)
      \item Align mitigations with the layer that fails (algorithm, implementation, network, operations)
    \end{itemize}
  \end{block}
\end{frame}

\begin{frame}
  \begin{block}{Takeaways}
  \begin{description}
    \item[Baseline] Proofs and formal methods give precise guarantees under assumptions
    \item[Bridging] Protocol and implementation verification reduce gaps to reality
    \item[Modeling] Choice of model (isolating vs mirror) shapes conclusions and validity
    \item[Practice] Report assumptions, base rates, and limitations; align methods to claims
  \end{description}
  \end{block}
\end{frame}


\section{Related to TAMOS}

\begin{frame}
  \begin{question}[Models]
    \begin{enumerate}
      \item How do we use models in security?
      \item How does that affect our results?
        \begin{itemize}
          \item Is ElGamal encryption secure or not?
        \end{itemize}
      \item What's the best we can do about the situation?
    \end{enumerate}
  \end{question}
\end{frame}

\begin{frame}
  \begin{question}[Engineering design]
    \begin{enumerate}
      \item Is our use of models just natural to the field or due to security 
        being an engineering discipline?
    \end{enumerate}
  \end{question}
\end{frame}

\section{Model checking and theorem proving}

\begin{frame}
  \begin{question}
    \begin{enumerate}
      \item How do model checking and theorem proving relate to each other and 
        cryptography?
      \item How sure can we be of the results when using model checking or 
        theorem proving?
        Better than the ElGamal example?
    \end{enumerate}
  \end{question}
\end{frame}
